% Nouveau EGA V (cote 25): archive cote 25, site page 79 (PDF page 80).
% typescript. Draft compared: none.
% AI-transcribed by Claude Opus 5.5 (claude-opus-5-5, adaptive thinking, effort high), Message Batches API, batch ega25b (msgbatch_01A7n5YuXubUPNK1zsepi8RG), 2026-11-13; not yet checked by eye.
\scanpage{25/79}
\mspage{\hand{74}}

Reste à prouver le lemme 13.3.  Comme d'habitude, il suffit de prouver
que  la sous-variété linéaire de \inserted{co}dimension n-1  passant par x générique
L\unsure{,}  possède les propriétés  a), b) , c) .  Pour  a), \inserted{ainsi que le contenu dimensionnel de b) \unsure{c)} (pour toute immersion projective donnée)} cela résulte
\hand{[}aussitôt de 2.3. (revu et corrigé dans le nº 12)  appliqué (comme dans un
raisonnement déjà fait dans le nº8) à l'espace projectif  des droites
passant par x, et l'image de Z dans ledit par projection conique à
partir de x. \hand{[}Il pourrait être utile d'ailleurs d'expliciter certains
\margin{|} résultats obttenus par cette méthode, concernant les sections linéaires
par des sous-variétés linéaires assujetties à passer par une sous-variété
linéaire fixe,  au cours du texte ou dans\hand{/}un nº à part\hand{]}.  Pour \inserted{la lissité dans} b) et c),
\xout{remplaçant} on peut grâce à a)  remplacer X et Y  resp. par X-Z et Y-Z
qui sont lisses, et on est  ramené à prouver ceci:  Soit f: X$\rightarrow$ P un
morphisme non ramifié,\hand{(}avec X lisse sur k, et soit  x $\in$P(k),\inserted{tel que x n'appartienne à l'image d'aucune composante de X de dim $<$ m ,} alors si
\hand{$\eta$} est le point générique  d\unsure{e} \xout{xxxxx} la sous-grassmanienne  de Grass$^{m}$(P)
formée des variétés linéaires L de codimension m passant par \xout{unexvariété} \inserted{x}
\xout{linéairexdonnéexdexcodimensionxm+1,} \inserted{\unsure{\&}} $\underline{X}^{(m)}_{\hand{\eta}}$ est  lisse sur k, tout au
\margin{(} moins si k est de car 0, et dans le cas contraire, à condition de rempla-%
cer f par $\phi_{n}f$ , n entier $\geq$ 2 quelconque.  C'est là  un remords au nº9,
qui résulte lui-même du remords suivant au nº 8:  Avec les notations de
8.8. , (supposant donc X irréductible, ce qui est loisible pour le Pb
que nous traitons), si  l'on a \inserted{co}dim T $\geq$ \deleted{n-}2  ou si  $\underline{Y}^{\text{sing}}\rightarrow$ T  est
génériquement étale (condition automatiquement vérifiée si k est de car 0,
ou à condition de remplacer f par $\phi_{n}f$ avec n$\geq$2, cf nº 9), alors pour
l'hyperplan \xout{générique} H\hand{$_{\eta}$} passant par x  générique, \deleted{\unsure{on a}} $\underline{X}^{(1)}_{\hand{\eta}}$ est lisse
de dimension n-1,  sauf dans le cas  où  on aurait  f(X) =\hand{$\{$}x\hand{$\}$}  (donc n=0)\hand{.}
Ce résultat admis,\hand{(}qui liquide évidemment le cas particulier m=1 de
notre remords, on obtient aussitôt le cas  m général  par récurrence sur

\margin{\unsure{75}}
